% =====================================================================
%  ECRformer-Lite / SEN12MS-CR cloud-removal pipeline
%  Complete technical record of the work performed up to 2026-10-15
%  Author of record: Vaishnav (yovyshh)
%  Compile with: pdflatex ecrformer_pipeline.tex   (run twice for refs)
% =====================================================================
\documentclass[10pt,a4paper]{article}

\usepackage[margin=1in]{geometry}
\usepackage{amsmath,amssymb}
\usepackage{array}
\usepackage{float}
\usepackage{xcolor}
\usepackage{listings}
\usepackage[hidelinks]{hyperref}

\lstset{
  basicstyle=\ttfamily\footnotesize,
  breaklines=true,
  columns=fullflexible,
  frame=single,
  framerule=0.3pt,
  rulecolor=\color{gray!50},
  backgroundcolor=\color{gray!8},
  showstringspaces=false,
  keywordstyle=\color{blue!70!black},
  commentstyle=\color{green!40!black},
  language=bash
}

\newcommand{\code}[1]{\texttt{#1}}

\title{\textbf{ECRformer-Lite for Sentinel-2 Cloud Removal}\\[0.3em]
\large Complete Pipeline, Experiments and Diagnostics\\[0.2em]
\normalsize Record of work to 2026-10-15}
\author{Vaishnav (yovyshh)}
\date{2026-10-15}

\begin{document}
\maketitle

\begin{abstract}
\noindent
This document records the full pipeline built on top of the ECRformer architecture for
multimodal optical--SAR cloud removal on a cached Sentinel-2 / Sentinel-1 dataset, together with
every experiment and diagnostic run performed so far. It states the data contract, the
architecture exactly as instantiated, the objective function, the optimisation and validation
protocol, the preserved checkpoint lineage, the chronological diagnostic workstreams, and the
currently running gradient-preservation pilot. A final section records the measured limitations
of the current baseline without softening them. Every numeric value quoted here was produced by a
script in the repository; values that are not recorded in the preserved artefacts are marked as
such rather than guessed.
\end{abstract}

\tableofcontents

% ---------------------------------------------------------------------
\section{Scope, task and notation}
% ---------------------------------------------------------------------

The task is \emph{cloud removal}: given a cloudy multi-band optical image and a co-registered
SAR image, produce a cloud-free optical image.

\begin{center}
\begin{tabular}{@{}ll@{}}
\hline
Symbol & Meaning \\
\hline
$S2_c$ & cloudy optical observation, 4 bands (B04, B03, B02, B08), reflectance in $[0,1]$ \\
$S1$   & Sentinel-1 SAR, 2 bands (VV, VH), stored in dB \\
$M$    & soft cloud mask, 1 band, $[0,1]$, cloud-positive \\
$Y$    & cloud-free optical reference (target) \\
$P$    & raw network prediction, \code{pred\_direct}, 4 bands, absolute reflectance \\
$\hat Y$ & mask-blended network output \\
\hline
\end{tabular}
\end{center}

The model input is the 7-channel tensor $X = [\,S1,\; S2_c,\; M\,]$ and the blended output is

\begin{equation}
\hat Y = (1-M)\odot S2_c \;+\; M \odot P .
\label{eq:blend}
\end{equation}

Equation~\eqref{eq:blend} is a \emph{policy} statement: clear pixels are returned verbatim from the
input. It is not an accuracy statement, and Section~\ref{sec:limits} quantifies what it costs.

% ---------------------------------------------------------------------
\section{Data}
% ---------------------------------------------------------------------

\subsection{Source and cache}

The upstream dataset is SEN12MS-CR / SEN12MS-CR-TS: paired Sentinel-1 and Sentinel-2 patches over
ROI tiles, each cloudy observation paired with a cloud-free observation of the same tile acquired
on a \emph{different date}. A pre-decoded, pre-aligned cache is used:

\begin{center}
\begin{tabular}{@{}ll@{}}
\hline
Item & Value \\
\hline
Cache root & \code{C:/Users/Windows 11 Pro/Videos/Projects/GAN/cache\_fall\_v3} \\
Training patches & 33{,}468 \\
Validation patches & 1{,}115 \\
Test split & held-out scene \code{s2\_139} --- \textbf{never accessed} \\
Patch size & $256\times256$ \\
Optical bands & B04 (red), B03 (green), B02 (blue), B08 (NIR) \\
File format & one \code{.npz} per patch per split \\
\hline
\end{tabular}
\end{center}

The cache lives in a sibling project directory (\code{Projects/GAN}), not inside the model
repository. This is a known coupling in the setup (Section~\ref{sec:limits}).

\subsection{Per-sample contract}

\code{data/cache\_dataset.py} returns a 5-tuple:

\begin{center}
\begin{tabular}{@{}lll@{}}
\hline
Index & Tensor & Shape \\
\hline
0 & \code{s2\_cloudy} & $(4,256,256)$ \\
1 & \code{s1}        & $(2,256,256)$ \\
2 & \code{s2\_clear} & $(4,256,256)$ \\
3 & \code{cf}        & scalar, stored cloud fraction \\
4 & \code{mask}      & $(1,256,256)$ \\
\hline
\end{tabular}
\end{center}

\subsection{SAR normalisation (on the fly, at load time)}

Four modes exist; the pipeline uses \code{paper\_exact}:

\begin{align*}
\text{VV}:&\quad \mathrm{clamp}(x,\,-25,\,0)/25 + 0 \quad\text{ i.e. } [-25,0]\,\text{dB} \to [0,1] \\
\text{VH}:&\quad \mathrm{clamp}(x,\,-35,\,0)/35 \qquad\text{ i.e. } [-35,0]\,\text{dB} \to [0,1]
\end{align*}

The alternatives \code{paper\_default} (both bands $[-25,0]$), \code{sen12ms\_resnet}
(VV $[-25,0]$, VH $[-32.5,0]$) and \code{raw} (unmodified dB) are retained and were
cross-checked in \code{tests/test\_final\_checks.py}.

\subsection{Mask refinement (three steps, applied every load)}

\begin{enumerate}
\item $9\times9$ mean filter, threshold at $0.5$ --- removes speckle:
      $m_1 = \mathbb{1}[\mathrm{avgpool}_{9\times9}(M) > 0.5]$
\item morphological dilation, radius 2 ($5\times5$ max-pool) --- grows the cloud region:
      $m_2 = \mathrm{maxpool}_{5\times5}(m_1)$
\item Gaussian blur ($k=11$, $\sigma=2$) --- produces the soft boundary:
      $M' = \mathrm{clamp}(\mathrm{gauss}_{11,2}(m_2), 0, 1)$
\end{enumerate}

All four tensors are cropped with the \emph{same} window, keeping pixel alignment.

\subsection{Cropping policy}

\begin{center}
\begin{tabular}{@{}ll@{}}
\hline
Split & Crop \\
\hline
Training & \code{crop\_size=128}, window sampled randomly per sample \\
Validation & \code{crop\_size=128}, window fixed to the image centre (deterministic) \\
Full-resolution diagnostics & \code{crop\_size=None} $\Rightarrow$ the whole $256\times256$ patch \\
\hline
\end{tabular}
\end{center}

The two evaluation regimes (masked 128 centre-crop vs.\ all-pixel full 256) are \textbf{not
comparable}; differences reach 2.3~dB with mixed sign. See Section~\ref{sec:limits}.

% ---------------------------------------------------------------------
\section{Architecture}
% ---------------------------------------------------------------------

\subsection{Wrapper: \code{ECRformerAdaptedModel}}

A thin adapter around the upstream backbone, adding decoupled stems, a 4-channel output head and
the multi-scale auxiliary heads needed by the SDFL loss.

\begin{center}
\begin{tabular}{@{}ll@{}}
\hline
Component & Configuration \\
\hline
Branch 0 (radar stem) & \code{Conv2d(2, 16, k=7, p=3)} on SAR \\
Branch 1 (optical stem) & \code{Conv2d(5, 16, k=7, p=3)} on $[S2_c, M]$ \\
Backbone & \code{ECRformerModel}, decoupled input \\
  & \code{num\_layers=4}, \code{num\_blocks=[2,2,1,1]} \\
  & \code{features\_start=32} (``Light'' width) \\
  & \code{bottle\_neck='tsa'}, \code{cbam='1ca2+1sa2'} \\
  & \code{num\_refine=0} (refinement path bypassed) \\
  & \code{conv\_type='conv'}, \code{norm\_type='batch'} \\
Output head & \code{final\_conv: Conv2d(64, 4, k=3, p=1)} \\
Auxiliary heads & 3 \code{down\_proj} (encoder) + 3 \code{up\_proj} (decoder), 4 ch each \\
Head init & \code{head\_init='mean\_prior'} \\
\hline
\end{tabular}
\end{center}

The \code{mean\_prior} initialisation sets the bias of \code{final\_conv} and of every auxiliary
head to the empirical clear-terrain band means
\code{CLEAR\_CHANNEL\_MEANS = [0.0847, 0.0919, 0.1127, 0.1884]} for B04/B03/B02/B08, and initialises
the corresponding weights from $\mathcal{N}(0, 0.005^2)$. The network therefore starts as a
constant clear-scene image and learns only a correction to it. This is deliberate --- it makes the
learning problem well-conditioned --- but it also biases the solution toward the mean
(Section~\ref{sec:limits}).

\subsection{Forward pass}

\begin{lstlisting}[language=Python]
x_in = cat([s1, s2_cloudy, mask], dim=1)          # [B, 7, H, W]
pred_direct, (down_projs, up_projs) = backbone(x_in)   # P, [B,4,H,W]
out = (1.0 - mask) * s2_cloudy + mask * pred_direct    # Eq. (1)
\end{lstlisting}

No hard-coded spatial constants exist anywhere in the backbone: all shapes are derived from the
input, so $256\times256$ inference needs no crop, resize or padding. (This was verified explicitly
in \code{diagnostics\_full256\_epoch30/REPORT.md}.)

\subsection{Repositories of record}

\begin{center}
\begin{tabular}{@{}ll@{}}
\hline
Role & Source \\
\hline
Upstream method & \code{github.com/zzaiyan/ECRformer} (ISPRS J.~P.~RS, 2026) \\
  & XCA + MDWA attention; SDFL training strategy \\
Task reference for video segmentation & \code{github.com/AMVerge-team/AMVerge} \\
\hline
\end{tabular}
\end{center}

% ---------------------------------------------------------------------
\section{Objective function}
% ---------------------------------------------------------------------

\subsection{SDFL loss (\code{losses/sdfl\_loss.py})}

Semantic-Decoupled Feature Learning supervises the encoder multi-scale features with a
structure-oriented criterion and the decoder features with a texture/radiometric criterion.

\paragraph{Encoder (down) branch --- structure-oriented, defaults $0.9/0.1$}
\begin{equation}
\mathcal{L}_{enc} = \frac{1}{N}\sum_{s=1}^{N}
\Big[ w^{enc}_{ssim}\big(1-\mathrm{SSIM}(D_s, Y_s)\big) + w^{enc}_{l1}\,\|D_s - Y_s\|_1 \Big]
\end{equation}

\paragraph{Decoder (up) branch --- texture-oriented, defaults $0.9/0.1$}
\begin{equation}
\mathcal{L}_{dec} = \frac{1}{N}\sum_{s=1}^{N}
\Big[ w^{dec}_{l1}\,\|U_s - Y_s\|_1 + w^{dec}_{ssim}\big(1-\mathrm{SSIM}(U_s, Y_s)\big) \Big]
\end{equation}

where $D_s, U_s$ are the \code{down\_proj} / \code{up\_proj} features at scale $s$, $Y_s$ is the
correspondingly downsampled target ($N=3$ scales, $Y_1 = Y$), and $\mathrm{SSIM}$ is computed by
\code{util/pytorch\_ssim} with an $11\times11$ window.

\paragraph{Main term --- full-domain supervision on the raw prediction}
With \code{full\_domain\_main=True} the main loss is applied to $P$ (not to $\hat Y$) and is
unmasked, so the network is penalised everywhere, including clear ground:
\begin{equation}
\mathcal{L}_{main} = w^{main}_{l1}\,\|P - Y\|_1
      + w^{main}_{ssim}\big(1 - \mathrm{SSIM}(P, Y)\big)
\end{equation}

A masked, area-normalised L1 on the blended output is still computed but only as a diagnostic:
\begin{equation}
\mathcal{L}^{diag}_{l1,cloudy} = \frac{\sum |\hat Y - Y| \odot M}{C\cdot\sum M + \varepsilon},
\qquad \varepsilon = 10^{-8},\; C = 4 .
\end{equation}

\paragraph{Total}
\begin{equation}
\mathcal{L}_{total} = (1-\lambda_{enc}-\lambda_{dec})\,\mathcal{L}_{main}
      + \lambda_{enc}\mathcal{L}_{enc} + \lambda_{dec}\mathcal{L}_{dec}
\end{equation}

The alternative \code{paper\_additive} form uses coefficient $1.0$ on $\mathcal{L}_{main}$ instead
of $0.9$. Both forms are asserted against hand-computed values to $<10^{-5}$ in the test suite.

\subsection{Weight selection actually used}

\begin{center}
\begin{tabular}{@{}lll@{}}
\hline
Parameter & Class default & Value used in the pilot \\
\hline
\code{total\_loss\_form} & \code{repo\_convex} & \code{repo\_convex} \\
\code{sdfl\_mode} & \code{official\_repo} & \code{official\_repo} (enc 0.9/0.1, dec 0.9/0.1) \\
\code{lambda\_enc} & 0.05 & 0.05 \\
\code{lambda\_dec} & 0.05 & 0.05 \\
\code{w\_main\_l1} & 0.9 & 1.0 \\
\code{w\_main\_ssim} & 0.1 & 0.0 \\
\hline
\end{tabular}
\end{center}

The weights used for the preserved epoch-30 baseline run are \textbf{not recorded in the
checkpoint} (\code{best\_model.pth} stores only the state dicts and a blended PSNR). This is a
documented gap, not an assumption.

\subsection{Gradient-preservation term (new, current experiment)}
\label{sec:gradloss}

An added finite-difference term, motivated by the measured high-frequency loss in
Section~\ref{sec:limits}:
\begin{align}
D_x(Q) &= Q[\dots,:,1:] - Q[\dots,:,:-1] \\
D_y(Q) &= Q[\dots,1:,:] - Q[\dots,:-1,:] \\
\mathcal{L}_{grad} &= \overline{|D_x(P)-D_x(Y)|} + \overline{|D_y(P)-D_y(Y)|} \\
\mathcal{L}_{total}^{new} &= \mathcal{L}_{total} + \lambda_{grad}\mathcal{L}_{grad}
\end{align}

Applied to the raw prediction $P$, all four bands, reflectance units, no resize or normalisation.

The weight was \emph{measured}, not chosen. On eight fixed-seed training-only batches (batch
size 4) with the epoch-30 EMA in \code{torch.inference\_mode()} and no backward pass:
\begin{equation}
\lambda_{grad} = \frac{0.05}{\operatorname{median}\!\big(\mathcal{L}_{grad}/\mathcal{L}_{others}\big)}
= 0.07943353602356962
\label{eq:lambda}
\end{equation}

\begin{center}
\begin{tabular}{@{}rrrrrr@{}}
\hline
Batch & $\mathcal{L}_{main}$ & $\mathcal{L}_{enc}$ & $\mathcal{L}_{dec}$
     & $\mathcal{L}_{others}$ & $\mathcal{L}_{grad}$ \\
\hline
0 & 0.008887 & 0.085063 & 0.016286 & 0.013065 & 0.005254 \\
1 & 0.014498 & 0.106879 & 0.021368 & 0.019461 & 0.011182 \\
2 & 0.011357 & 0.090140 & 0.017162 & 0.015586 & 0.009424 \\
3 & 0.008130 & 0.056070 & 0.011343 & 0.010688 & 0.006526 \\
4 & 0.012301 & 0.135141 & 0.020505 & 0.018853 & 0.012224 \\
5 & 0.008031 & 0.053215 & 0.011741 & 0.010475 & 0.007184 \\
6 & 0.008612 & 0.059500 & 0.012609 & 0.011357 & 0.009234 \\
7 & 0.011279 & 0.075925 & 0.016575 & 0.014776 & 0.009719 \\
\hline
\end{tabular}
\end{center}

Median weighted share $5.00\%$, maximum $6.46\%$ (guard $\le 10\%$). Artefact:
\code{experiments/ecrformer\_gradient\_pilot/probe\_loss\_scale.json}.

\emph{Correction of record.} An earlier plan draft stated
$\lambda = 0.05\cdot\operatorname{median}(\mathcal{L}_{others}/\mathcal{L}_{grad})$. That form does
not generally place the median weighted contribution at 5\%. The implemented form,
Eq.~\eqref{eq:lambda}, does, and the table above verifies it.

% ---------------------------------------------------------------------
\section{Optimisation and training loop}
% ---------------------------------------------------------------------

\begin{center}
\begin{tabular}{@{}ll@{}}
\hline
Item & Setting \\
\hline
Optimiser & Adam (PyTorch defaults) \\
Gradient clipping & \code{clip\_grad\_norm\_(max\_norm=1.0)} per optimiser update \\
Gradient accumulation & \code{accum\_steps} micro-batches per update \\
LR schedule & \code{CosineAnnealingLR(T\_max = epochs, eta\_min = min\_lr)} \\
EMA & \code{update\_ema(ema\_model, model, decay=0.999)} on every update \\
Precision & \code{fp32} (default) or \code{bf16} under \code{torch.autocast} \\
DataLoader workers & configurable (\code{--num\_workers}) \\
Checkpoint selection & \code{--select\_metric raw|blend} \\
\hline
\end{tabular}
\end{center}

EMA is applied parameter-wise over the floating-point entries of both state dicts and the
\code{ema\_state\_dict} is what is used for scoring.

\subsection{Validation protocol}

\code{validate()} evaluates on cloudy pixels (\code{mask > 0.1}) and stratifies into six buckets:
\code{all}, \code{light}, \code{medium}, \code{heavy}, \code{heavy\_opaque}, \code{heavy\_hazy}.
Two tensors are scored from the \emph{same} forward pass:

\begin{itemize}
\item \code{raw} $=P$ (\code{pred\_direct}) --- the primary experimental result;
\item \code{blend} $=\hat Y$ --- comparison only; in clear regions it is identical to the input.
\end{itemize}

Per-variant metrics are \code{psnr}, \code{ssim}, \code{l1}, \code{pearson}, \code{ndvi\_mae},
\code{ndvi\_r} plus per-band L1/means. The legacy unsuffixed keys alias whichever variant
\code{--select\_metric} names. Each epoch also records band statistics of prediction and target
so that drift toward the mean is visible during training.

% ---------------------------------------------------------------------
\section{Checkpoint and run lineage}
% ---------------------------------------------------------------------

\subsection{Preserved baseline}

\begin{center}
\begin{tabular}{@{}ll@{}}
\hline
Item & Value \\
\hline
Path & \code{checkpoints/ecrformer\_light\_128\_50ep/best\_model.pth} \\
Epoch & 30 \\
Weights scored & \code{ema\_state\_dict} \\
Stored \code{psnr} & 30.1982 \emph{(blended semantics)} \\
\code{psnr\_select} key & \textbf{absent} --- selection metric not recorded \\
\code{scheduler\_state\_dict} & \textbf{absent} --- resume position unknowable \\
File size & 58{,}406{,}779 bytes \\
SHA-256 & \code{2f02aca75619a8b4db52ca5ef4910a0ad305e80c68d48c035c0bb975e5b1a5b6} \\
\hline
\end{tabular}
\end{center}

\code{latest\_model.pth} in the same directory is at epoch 32 with optimiser LR
$2.2222\times10^{-5}$ (versus $1.4474\times10^{-4}$ in the best checkpoint): training continued two
epochs past the best point without improving the selection metric, and no record explains why it
stopped there.

\subsection{Archive for reversion}

The baseline checkpoint, twelve full-resolution output pairs, comparison figures, metric tables
and a README were published as a separate private repository,
\code{github.com/yovyshh/ecrformer-lite-baseline-archive}, precisely so that the epoch-30 state can
be recovered independently of the working tree. The checkpoint digest is recorded in
\code{checkpoint/SHA256SUMS} inside that repository and matches the table above. The archive
deliberately excludes the dataset cache, the held-out scene, training code and later checkpoints.

\subsection{All run directories present}

\begin{center}
\begin{tabular}{@{}ll@{}}
\hline
Run & Note \\
\hline
\code{ecrformer\_light\_128\_stage1} & stage-1 run \\
\code{ecrformer\_light\_128\_50ep} & \textbf{preserved baseline} (best = epoch 30) \\
\code{ecrformer\_light\_sdfl} & SDFL integration run \\
\code{ecrformer\_light\_sdfl\_run1\_diverged} & \textbf{diverged} --- kept as evidence \\
\code{ecrformer\_light\_sdfl\_run2..run5} & iterations \\
\code{ecrformer\_light\_sdfl\_run6} & last pre-cache run \\
\code{experiments/ecrformer\_gradient\_pilot} & current, isolated (Section~\ref{sec:pilot}) \\
\hline
\end{tabular}
\end{center}

Upstream pretrained weights are kept under \code{pretrained/} (\code{full.ckpt},
\code{light.ckpt}) and are not modified.

% ---------------------------------------------------------------------
\section{Diagnostic workstreams (chronological)}
% ---------------------------------------------------------------------

All directories below were produced read-only: no checkpoint, cache, model or history file was
modified by any of them, and \code{s2\_139} was never opened.

\begin{center}
\begin{tabular}{@{}llp{8.2cm}@{}}
\hline
Time (2026) & Directory & Purpose / finding \\
\hline
13:22 & \code{diagnostics\_epoch30} & epoch-30 reference outputs \\
13:50 & \code{diagnostics\_cloud\_coverage} & behaviour stratified by cloud coverage \\
14:09 & \code{diagnostics\_model\_quality} & output-quality inspection \\
14:29 & \code{diagnostics\_fixes} & resume / blend / scheduler fixes and their pre-verification \\
14:30 & \code{diagnostics\_raw\_outputs} & raw \code{pred\_direct} saved unclipped \\
14:41 & \code{diagnostics\_full256\_epoch30} & full $256\times256$ inference, 12 patches, 6-panel figures; established that resolution does not recover sharpness \\
15:00 & \code{diagnostics\_smoothing\_audit} & quantified the over-smoothing \\
15:11 & \code{diagnostics\_regime\_audit} & regime stratification audit \\
15:25 & \code{diagnostics\_regime\_27x27} & small-patch regime study \\
15:39 & \code{diagnostics\_scene\_compare} & cross-scene comparison \\
17:55 & \code{diagnostics\_shadow\_audit} & shadow / dark-surface hypothesis ruled out or bounded \\
18:09 & \code{diagnostics\_ocm\_stage0} & occlusion-mask stage-0 study \\
18:28 & \code{diagnostics\_ocm\_stage1} & occlusion-mask stage-1 study \\
19:03 & \code{diagnostics\_solution\_study} & 8-patch full-resolution baseline metric table \\
19:16 & \code{diagnostics\_five\_coverage\_stack} & the corrected 5-panel coverage figure (verified) \\
19:43 & \code{diagnostics\_synthetic\_haze\_v2} & training-only synthetic-corruption feasibility \\
19:50 & \code{diagnostics\_project\_review} & consolidated read-only review \\
\hline
\end{tabular}
\end{center}

Two results from this phase matter downstream.

\paragraph{The corrected five-coverage figure is trustworthy.} Its five displayed PSNRs were
independently recomputed at full resolution with the epoch-30 EMA and matched to within 0.01~dB;
the mask polarity (cloud-positive) and Eq.~\eqref{eq:blend} were both confirmed:

\begin{center}
\begin{tabular}{@{}lrrrr@{}}
\hline
Sample & Coverage & Input & Raw & Blend \\
\hline
\code{s2\_65\_p88} & 2.44\% & 35.37 & 37.40 & 35.38 \\
\code{s2\_3\_p632} & 12.55\% & 23.98 & 32.62 & 29.85 \\
\code{s2\_40\_p476} & 34.99\% & 14.45 & 28.42 & 27.33 \\
\code{s2\_40\_p368} & 62.47\% & 11.66 & 34.57 & 27.31 \\
\code{s2\_3\_p680} & 87.26\% & 10.88 & 29.14 & 28.20 \\
\hline
\end{tabular}
\end{center}

\paragraph{Synthetic haze is not yet usable.} Ninety-six deterministic training recipient/donor
pairs were generated. The donor-mask veil matched the failure PSNRs but its standard-deviation
ratios (2.88/4.20/5.67/1.52) fell short of the real failures (3.60--6.27) and its residual-gradient
ratios were \emph{lower} in every band; the donor-residual overlay produced standard-deviation
ratios near 8--15 and 77 of 96 candidates out of $[0,1]$. Neither is physically admissible for
training, so no augmentation pilot was launched.

% ---------------------------------------------------------------------
\section{The gradient-preservation pilot (running)}
\label{sec:pilot}
% ---------------------------------------------------------------------

Everything in this section is confined to \code{experiments/ecrformer\_gradient\_pilot/}. The
baseline directories are never written to.

\subsection{Pre-flight and safety}

\begin{itemize}
\item Pre-change \code{git status} recorded to \code{backups/pre\_change/preflight.txt}. Most
project sources are untracked, including \code{train\_ecrformer.py} itself, so Git cannot restore
them; a byte copy was therefore taken to
\code{backups/pre\_change/train\_ecrformer.py}. \code{git reset --hard}, \code{git clean},
\code{git checkout --} and \code{git restore} were never run.
\item New CLI flags: \code{--lambda\_gradient}, \code{--init\_from\_ema}, \code{--experiment\_dir}.
A positive \code{--lambda\_gradient} is rejected unless it is paired with \code{--init\_from\_ema}
and \code{--experiment\_dir}, and run names containing path traversal are rejected.
\item A fresh run is initialised from the epoch-30 EMA with a \emph{fresh} optimiser and scheduler:
no optimiser or scheduler state is loaded, and \code{--resume} is never used.
\end{itemize}

\subsection{Test-first evidence}

\begin{lstlisting}
python -m pytest tests/test_ecrformer_gradient_pilot.py \
  tests/test_resume_reference_schedule.py \
  tests/test_blend_and_scheduler_fixes.py \
  tests/test_final_checks.py -q
# -> 35 passed
\end{lstlisting}

RED was observed first (missing module, then a genuine \code{IndexError} in the gradient helper
for 3-dimensional tensors). Coverage includes: zero weight leaves \code{loss\_total} bit-identical
to the existing criterion and preserves \code{loss\_main}/\code{loss\_enc\_aux}/\code{loss\_dec\_aux};
non-zero weight adds exactly $\lambda\mathcal{L}_{grad}$; the EMA loader consumes
\code{ema\_state\_dict} and leaves optimiser/scheduler fresh; unsafe option combinations are
rejected.

Two pre-existing defects were fixed in the course of this: \code{tests/test\_final\_checks.py}
still constructed the model with removed keywords (\code{residual\_blend}, \code{zero\_init\_final})
and asserted a residual-additive blend that the code no longer implements; and a stray \code{\%}
in the \code{--sdfl\_mode} help string made \code{--help} raise
\code{ValueError: unsupported format character}.

\subsection{Baseline reference (8 fixed validation patches, full $256\times256$)}

\begin{center}
\begin{tabular}{@{}lrrrrrr@{}}
\hline
Sample & Mask \% & Raw PSNR & Blend PSNR & Raw SSIM & grad-L1 & grad-energy \\
\hline
\code{s2\_119\_p208} & 0.0 & 24.94 & 10.76 & 0.8000 & 0.014969 & 0.322 \\
\code{s2\_119\_p236} & 0.9 & 22.79 & 10.63 & 0.7424 & 0.019154 & 0.263 \\
\code{s2\_119\_p92}  & 0.0 & 27.87 & 13.61 & 0.8668 & 0.011594 & 0.477 \\
\code{s2\_40\_p268}  & 5.6 & 28.20 & 28.37 & 0.9411 & 0.010552 & 0.639 \\
\code{s2\_65\_p328}  & 0.0 & 27.46 & 31.14 & 0.9358 & 0.011697 & 0.768 \\
\code{s2\_65\_p124}  & 0.8 & 42.85 & 41.39 & 0.9842 & 0.004855 & 0.846 \\
\code{s2\_119\_p100} & 100.0 & 26.51 & 26.51 & 0.9037 & 0.010812 & 0.640 \\
\code{s2\_65\_p336}  & 89.3 & 34.43 & 34.25 & 0.9542 & 0.007739 & 0.902 \\
\hline
\textbf{mean} & & \textbf{29.38} & \textbf{24.58} & \textbf{0.8910} & \textbf{0.011421} & \textbf{0.607} \\
\hline
\end{tabular}
\end{center}

Artefacts: \code{baseline\_reference/per\_patch\_metrics.csv}, \code{manifest.json},
\code{predictions/*.npz}.

\subsection{Setup gate}

\code{verify\_setup.py} $\rightarrow$ \code{PASS\_SETUP\_GATE}:
\code{baseline\_epoch=30}, \code{ema\_tensor\_match=True}, \code{fresh\_ema\_match=True},
\code{fresh\_optimizer=True}, \code{fresh\_scheduler=True},
\code{experiment\_paths\_isolated=True}, \code{test\_scene\_accessed=False},
smoke pass finite and inference-only.

\subsection{Run configuration}

\begin{lstlisting}
python train_ecrformer.py \
  --init_from_ema checkpoints/ecrformer_light_128_50ep/best_model.pth \
  --experiment_dir experiments/ecrformer_gradient_pilot \
  --run_name gradloss_3ep \
  --lambda_gradient 0.07943353602356962 \
  --epochs 3 --micro_batch_size 4 --accum_steps 4 \
  --lr 1.447417e-4 --min_lr 1e-5 --crop_size 128 \
  --sar_norm paper_exact --sdfl_mode official_repo \
  --total_loss_form repo_convex \
  --lambda_enc 0.05 --lambda_dec 0.05 --w_main_l1 1.0 --w_main_ssim 0.0 \
  --num_refine 0 --precision fp32 --num_workers 0 --select_metric raw
\end{lstlisting}

Effective batch size 16; fresh \code{CosineAnnealingLR(T\_max=3, eta\_min=1e-5)}; started at the
documented epoch-30 effective LR $1.447417\times10^{-4}$.

\subsection{Live status (2026-10-15)}

\begin{center}
\begin{tabular}{@{}ll@{}}
\hline
Item & Value \\
\hline
Launched & 08:53 IST, detached (\code{run\_gradloss\_3ep.bat}) \\
Epoch 1 & complete 09:29; LR $1.447417\times10^{-4}$ for all 389 recorded updates \\
Epoch 2 & in progress; LR stepped to $1.110563\times10^{-4}$ as the cosine schedule requires \\
Epoch duration & $\approx 29$--36 min (33{,}468 patches, batch 4, crop 128) \\
GPU & RTX 5070, $\approx$7.6~GB of 12.2~GB resident \\
Checkpoints & \code{checkpoints/gradloss\_3ep/\{best,latest\}\_model.pth} (58.4~MB each) \\
LR trace & \code{outputs/gradloss\_3ep/optimizer\_lr\_trace.csv} \\
\hline
\end{tabular}
\end{center}

\subsection{Pre-declared success gates}

\begin{enumerate}
\item Both \code{s2\_119\_p208} and \code{s2\_119\_p236} improve target-relative gradient error
      \emph{and} move gradient-energy ratio toward 1.
\item None of the four critical samples (\code{p208}, \code{p236}, \code{p92}, \code{p40\_p268})
      loses more than $0.5$~dB raw PSNR.
\item Cohort mean raw PSNR regresses by no more than $0.1$~dB.
\item Proposed non-inferiority margins: cohort mean SSIM drop $\le 0.005$; cohort mean MAE rise
      $\le 0.001$; absolute per-band mean-bias increase $\le 0.005$; out-of-$[0,1]$ fraction does
      not rise by more than $0.1$ percentage point.
\item All losses, gradients and outputs finite; no invented edges or amplified noise in the saved
      comparisons.
\end{enumerate}

These thresholds are taken from the plan authored before the run, but the numeric margins in
gate 4 were never separately confirmed by the user, and \code{REPORT.md} was written immediately
\emph{after} launch. Both facts are recorded rather than glossed.

Failure of any gate means FAILED: keep the logs and metrics, do not promote the checkpoint, and
proceed to rollback verification. A pass is a screening result only --- it does not replace the
baseline and it claims no statistical significance.

An additional gate follows from Section~\ref{sec:limits}: if \code{p208}/\code{p236} are
mis-paired samples in the cache, gate 1 tests the cache rather than the loss. That pairing check
has not been performed.

% ---------------------------------------------------------------------
\section{Measured limitations of the current baseline}
\label{sec:limits}
% ---------------------------------------------------------------------

\subsection{Systematic over-smoothing}

Across twelve full $256\times256$ patches the raw prediction retains only
$23$--$68\%$ of the target's high-frequency energy (cohort mean $0.607$), the regression slope is
$0.48$--$0.72$, and per-band standard-deviation ratios fall as low as $0.156$ (red band,
\code{s2\_119\_p208}). On that patch the \emph{overall} standard-deviation ratio is $0.765$ while
the gradient ratio is $0.274$: the output has roughly the right low-frequency content and almost no
structure. Increasing the evaluation area from $128^2$ to $256^2$ bought no sharpness at all.

\subsection{The objective has no sharpness incentive}

L1 plus SSIM on absolute reflectance is minimised by a conditional mean. For a multi-modal problem
the conditional mean is blurry by construction, and the multi-scale auxiliary terms inherit the
same property. The \code{mean\_prior} head initialisation starts the network at a constant
clear-scene image, which reinforces the bias. Nothing in the original objective ever rewarded an
edge. This is the hypothesis the pilot tests and the only thing it changes.

\subsection{PSNR conceals the defect}

Cohort mean raw PSNR is $29.38$~dB. The most flattering numbers occur precisely where the output is
flattest: \code{p56} reports $44.09$~dB at $100\%$ coverage, and \code{p208} jumps from $10.76$ to
$24.94$~dB while sitting at the bottom of the sharpness ranking ($0.274$).

\subsection{The cloud mask is unreliable in both directions}

\begin{center}
\begin{tabular}{@{}lrr@{}}
\hline
Sample & Input PSNR & Mask coverage \\
\hline
\code{s2\_134\_p56}   & 8.40 & 100.0\% \\
\code{s2\_3\_p680}    & 10.88 & 87.3\% \\
\code{s2\_119\_p208}  & \textbf{10.76} & \textbf{0.0\%} \\
\code{s2\_119\_p236}  & \textbf{10.63} & \textbf{0.9\%} \\
\code{s2\_119\_p100}  & 23.45 & 100.0\% \\
\code{s2\_119\_p264}  & 23.03 & 5.4\% \\
\hline
\end{tabular}
\end{center}

\code{p208} is as degraded as \code{p56} and \code{p680} yet its mask declares the patch entirely
clear; \code{p100} is barely degraded and its mask claims total cover. The cache's own \code{cf}
field agrees with the mask ($0.0012$ for \code{p208}), so two independent cloud measures are
contradicted by the radiometry of the same patch. Consequences: Eq.~\eqref{eq:blend} returns the
cloudy input on \code{p208}/\code{p236}, discarding $12$--$14$~dB of the network's own improvement
(blend $10.76$ versus raw $24.94$), and the pilot's gate 1 is anchored on exactly these samples.

\subsection{Temporal confounding}

SEN12MS-CR-TS pairs a cloudy date with a \emph{different} clear date. A patch with an empty mask and
poor input PSNR may be poor because of seasonal or land-cover change rather than cloud, and the
network can gain PSNR there by predicting the terrain prior. That is a plausible reading of the
flat, low-gradient outputs on \code{p208}/\code{p236}. The evaluation does not currently separate
the two effects, which weakens any claim about cloud removal specifically.

\subsection{Two incompatible metric families}

Training validation is $128\times128$ centre-cropped, mask-restricted and (historically)
blend-selected. The diagnostics are $256\times256$, all-pixel and raw. Differences reach $2.3$~dB
with mixed sign, so the two families cannot be merged or directly compared. The preserved
checkpoint's stored \code{psnr} of $30.1982$ is blended and its selection key is missing, while the
current default is raw --- a silent change of meaning.

\subsection{State that cannot be resumed}

Neither preserved checkpoint stores \code{scheduler\_state\_dict}, so the schedule position is
unknowable; a plain resume keeps the saved LR, and a correct resume would need the original
\code{epochs}, \code{lr} and \code{min\_lr}, which were not recorded. \code{--resume} loads the
\emph{latest} checkpoint together with its optimiser state and therefore cannot start a fresh run
from the epoch-30 EMA --- the gap that motivated the new \code{--init\_from\_ema} path.

\subsection{Unusable synthetic augmentation}

See the end of Section~7: neither donor-mask veil nor donor-residual overlay reproduces the
statistics of the real failure patches, and the latter violates reflectance bounds on $80\%$ of
candidates.

\subsection{Engineering debt}

\begin{itemize}
\item Most sources, including \code{train\_ecrformer.py} and \code{data/cache\_dataset.py}, are
      untracked in Git; there is no repository-level rollback path, only manual backups.
\item The cache lives in a sibling project directory with no version pin.
\item The trainer prints with default buffering, so a redirected run log stays empty until the
      process exits; \code{python -u} should be used.
\item The historical $256\times256$ \emph{training} out-of-memory failure is real (activations,
      gradients and optimiser state); inference at the same size needs only 347~MiB.
\item No statistical significance is claimed anywhere; the 8-patch cohort is a screening set
      chosen to include known failures, and the only genuinely held-out scene is unmeasured.
\end{itemize}

% ---------------------------------------------------------------------
\section{Reproducibility commands}
% ---------------------------------------------------------------------

\begin{lstlisting}
cd "/c/Users/Windows 11 Pro/Videos/Projects/ECRformer"
PY="C:/Users/Windows 11 Pro/AppData/Local/Programs/Python/Python312/python.exe"

# focused test suite (expect 35 passed)
"$PY" -m pytest tests/test_ecrformer_gradient_pilot.py \
  tests/test_resume_reference_schedule.py \
  tests/test_blend_and_scheduler_fixes.py tests/test_final_checks.py -q

# measured lambda for the gradient term (train split only, inference mode)
"$PY" experiments/ecrformer_gradient_pilot/probe_loss_scale.py

# baseline reference on the 8 fixed validation patches
"$PY" experiments/ecrformer_gradient_pilot/evaluate_pilot.py --mode baseline

# setup gate before any training
"$PY" experiments/ecrformer_gradient_pilot/verify_setup.py

# score the experimental checkpoint against the baseline
"$PY" experiments/ecrformer_gradient_pilot/evaluate_pilot.py --mode experimental \
  --checkpoint experiments/ecrformer_gradient_pilot/checkpoints/gradloss_3ep/best_model.pth
\end{lstlisting}

Launcher for the detached run: \code{run\_gradloss\_3ep.bat}.

\subsection{Rollback}

\begin{lstlisting}
cd "/c/Users/Windows 11 Pro/Videos/Projects/ECRformer"
cp -p experiments/ecrformer_gradient_pilot/backups/pre_change/train_ecrformer.py train_ecrformer.py
sha256sum train_ecrformer.py \
  experiments/ecrformer_gradient_pilot/backups/pre_change/train_ecrformer.py
\end{lstlisting}

Optionally remove only the newly added code and test files; never the experiment directory, the
baseline reference, the report, the logs or the cache. If the pilot fails, move only
\code{checkpoints/gradloss\_3ep/} to \code{failed\_artifacts/}. Never delete the baseline or any
earlier diagnostic.

% ---------------------------------------------------------------------
\section{Artefact map}
% ---------------------------------------------------------------------

\begin{center}
\begin{tabular}{@{}lp{9.3cm}@{}}
\hline
Path & Contents \\
\hline
\code{data/cache\_dataset.py} & cached SEN12MS-CR loader, SAR normalisation, mask refinement, cropping \\
\code{models/ecrformer\_adapted.py} & task adapter: stems, heads, Eq.~\eqref{eq:blend} \\
\code{models/ecrformer\_model.py} & upstream ECRformer backbone (XCA, MDWA, TSA, CBAM) \\
\code{losses/sdfl\_loss.py} & SDFL objective of Section~4 \\
\code{train\_ecrformer.py} & training loop, validation, EMA, checkpointing, LR trace \\
\code{tests/} & scheduler-resume, blend, SAR-normalisation and pilot tests \\
\code{checkpoints/ecrformer\_light\_128\_50ep/} & \textbf{protected baseline} (epoch 30 EMA) \\
\code{pretrained/} & upstream \code{full.ckpt}, \code{light.ckpt} \\
\code{diagnostics\_*/} & seventeen read-only diagnostic workstreams (Section~7) \\
\code{experiments/ecrformer\_gradient\_pilot/} & current pilot: code, probes, gates, report, logs \\
\code{docs/ecrformer\_pipeline.tex} & this document \\
\hline
\end{tabular}
\end{center}

% ---------------------------------------------------------------------
\section{Open items}
% ---------------------------------------------------------------------

\begin{enumerate}
\item Score the pilot against the pre-declared gates once epochs 2 and 3 finish, and fill in the
      experimental results section of \code{REPORT.md}.
\item Verify whether \code{s2\_119\_p208} and \code{s2\_119\_p236} are correctly paired and
      masked in the cache against the upstream SEN12MS-CR-TS tiles. This decides whether the
      observed defects are data or model, and whether gate 1 is meaningful.
\item Decide the reporting convention: raw all-pixel at full resolution, or masked crop, or both
      stated explicitly --- the two families must never be merged.
\item Decide whether the mask defect is worth fixing before further training, since it currently
      determines whether the blended product is usable at all.
\item Record the selection metric and scheduler state in every future checkpoint so that a run is
      resumable without guesswork.
\end{enumerate}

\vfill
\noindent\rule{\textwidth}{0.3pt}\\
\footnotesize
Every numeric value in this document was produced by a script in the repository. Values absent
from the preserved artefacts --- the baseline run's loss weights, its selection metric, and the
scheduler position of the legacy checkpoints --- are marked as unrecorded rather than inferred.
The held-out scene \code{s2\_139} was not accessed at any point described here.

\end{document}
